\section{Related Work}
\label{sec:related_work}

\paragraph{Reasoning-time sampling, search, and self-consistency.}
Chain-of-thought prompting \citep{wei2022chain} and zero-shot chain-of-thought \citep{kojima2022large} showed that explicitly eliciting intermediate reasoning improves large-language-model performance on mathematical and symbolic tasks.
Self-consistency extends this idea by sampling multiple reasoning paths and returning the most frequent final answer \citep{wang2023self}.
Tree-structured reasoning methods, including Tree of Thoughts \citep{yao2023tree} and Graph of Thoughts \citep{besta2024graph}, further broaden inference-time exploration by organizing candidate reasoning paths into explicit search structures.
Program-aided language-model reasoning similarly generates executable or semi-executable solution traces before extracting a final answer \citep{gao2023pal}.
The present work is closest to this family in that it assumes multiple candidate answers and runtime traces are available, but it studies a different question: how to choose among already generated answers using gold-free execution metadata rather than selecting by raw frequency alone.

\paragraph{Budgeted inference and test-time compute.}
Recent work on test-time scaling studies how additional inference computation can improve reasoning accuracy \citep{snell2024scaling}.
Budget-forcing methods such as s1 shape the amount of reasoning performed during generation by extending or truncating traces toward a target budget \citep{muennighoff2025s1}, while broader analyses of inference-time interventions compare sampling, voting, search, and allocation strategies under cost constraints \citep{wu2024inference}.
Our experiments include prompt-budgeting and support-based answer-selection baselines as evaluated comparators; they are introduced in the experimental setup rather than treated as established prior work.
The contribution here is therefore not a new generation-time scaling law; it is a post-generation answer-selection rule evaluated after candidate-producing methods have run under the same per-method budget cap.

\paragraph{Verifier-guided reranking and answer selection.}
A complementary line of work trains verifiers or reward models to rerank candidate solutions.
For GSM8K, \citet{cobbe2021training} showed that an outcome verifier can substantially improve answer selection when trained on labeled solution data.
Process reward models provide finer-grained supervision by scoring intermediate reasoning steps rather than only final answers \citep{lightman2024lets}.
Lightweight supervised text-classification frameworks such as SetFit can also support task-specific reranking when suitable labels are available \citep{tunstall2022efficient}.
In contrast, the policy studied here does not train a verifier and does not use correctness labels at decision time; it relies on runtime-visible frontier metadata and agreement patterns among independently generated candidate answers.

\paragraph{Cost-aware routing, adaptive computation, and positioning.}
Cost-aware routing systems such as FrugalGPT \citep{chen2023frugalgpt} and language-model cascades \citep{dohan2022language} reduce inference cost by deciding which model or stage to invoke for a given input.
Adaptive computation methods make a related decision within a model or computation graph, allocating more processing to inputs that appear to require it \citep{graves2016adaptive}.
The setting in this paper is different: the evaluated candidate-producing methods are compared under a matched per-method budget, and FIX-2+FIX-4 itself is a deterministic selection layer that issues no additional model calls once candidate answers and metadata are available.
The resulting claim is intentionally narrow: on the tested Cohere/GSM8K setting, two runtime-legal signals---weak frontier traces and unanimous external disagreement---provide useful information for selecting among already generated answers.
